\problema{Validate Binary Search Tree}{98}{src/02-arboles/98-validate-binary-search-tree.cpp}{M}

Dada la raíz de un árbol binario, determina si es un árbol de búsqueda binaria
(BST) válido. Para que un BST sea válido, el subárbol izquierdo de cada nodo
debe contener solo nodos con claves \textbf{menores} que la clave del nodo, el
subárbol derecho debe contener solo nodos con claves \textbf{mayores}, y ambos
subárboles deben ser también BST válidos.

\paragraph{Intuición.} El error más común al resolver este problema es verificar
solamente que el hijo izquierdo sea menor y el hijo derecho sea mayor que el
nodo actual. Esto \emph{no} es suficiente. Un nodo que esté muy abajo en el
subárbol derecho de la raíz podría tener un valor menor que la raíz, violando
así la propiedad fundamental del BST (donde \textbf{todo} el subárbol derecho
debe ser mayor que la raíz).

\begin{center}
\ttfamily
\begin{tabular}{ccc}
 & 5 & \\
1 & & 4\\
 & 3 & 6\\
\end{tabular}
\end{center}

\noindent En este ejemplo, $3 < 4$ y $6 > 4$, por lo que el nodo 4 y sus hijos
parecen correctos localmente. Sin embargo, el nodo $3$ está en el subárbol
derecho de la raíz $5$, pero no es mayor que $5$. ¡El árbol es inválido!

\paragraph{Solución.} En lugar de verificar solo relaciones padre-hijo, debemos
mantener un \emph{rango válido} \texttt{(minimo, maximo)} a medida que bajamos
por el árbol en nuestra recursión. Todo nodo que visitemos debe caer
estrictamente dentro de ese rango.
\begin{itemize}
    \item Al principio, la raíz puede tener cualquier valor: \texttt{(LONG\_MIN, LONG\_MAX)}.
    \item Cuando bajamos al hijo izquierdo, el límite \textbf{superior} del rango se actualiza al valor del nodo actual.
    \item Cuando bajamos al hijo derecho, el límite \textbf{inferior} del rango se actualiza al valor del nodo actual.
\end{itemize}

Usamos el tipo \texttt{long} de C++ para los límites \texttt{minimo} y
\texttt{maximo} porque los valores de los nodos pueden ser \texttt{INT\_MIN} o
\texttt{INT\_MAX}, y si usáramos \texttt{int}, fallarían las comparaciones en
esos casos extremos (overflow/underflow). El código:

\lstinputlisting{src/02-arboles/98-validate-binary-search-tree.cpp}

\complejidad{$O(n)$}{$O(h)$}

\paragraph{Notas de entrevista (Amazon).} Cuidado especial con los casos borde
como los límites de los enteros de 32 bits (\texttt{INT\_MAX} y \texttt{INT\_MIN}).
Es preferible usar \texttt{long} o parámetros de puntero nulo para representar el
infinito. Además, asegúrate de mencionar que el espacio auxiliar de la recursión es
$O(h)$ donde $h$ es la altura del árbol, lo cual se convierte en $O(n)$ en el
peor caso (árbol degenerado como una lista ligada) y $O(\log n)$ en el mejor caso
(árbol balanceado).
